\section{Introduction}\label{sec:intro}
Maintainers ask what a change will affect~\cite{sillito2006questions,ko2007information}; LLM explanations answer ``what
does this code do?'' and helped developers finish tasks without measurably improving understanding~\cite{nam2024gilt}. Tracing and explaining are distinct skills~\cite{lister2009tracing,murphy2012eipe};
predicting how behaviour changes under an edit is the skill maintenance needs, and the ``modify'' step of
PRIMM~\cite{sentance2019primm} already trains it by hand.

We propose the \emph{counterfactual explanation} of code: one intervention, the set of inputs on which the changed
function behaves differently, \emph{established by execution}, and an articulation of that difference whose parts can be
\emph{checked by execution}: a witness input and an executable predicate for the affected inputs. The obvious attack is
that LLMs misdescribe behavioural differences~\cite{khatib2026mutation,reddy2026articulate}, so this paper first measures
how often LLM-written explanations are wrong, and what execution grounding changes.
\begin{figure}[!t]
\centering
% Figure 1: comic-style example (real data: HumanEval+/52 below_threshold, CMP_OP `<` -> `<=`).
% Sources: artifact/results/mutants.jsonl (program, mutant, differing inputs),
% probe_stripped_qwen2.5-coder_14b.jsonl (DESC and CF answers on [[1,3,4],4]),
% predicates_claude-haiku-5-5.jsonl / predicates_qwen2.5-coder_14b.jsonl (ungrounded predicates, exactness on T).
% Icons: Twemoji 14.0.2, CC-BY 4.0 (see figures/icons/SOURCES.md).
% Accent colours, one meaning each: accOp = the changed operator; accOut = the contested outcome (guess vs. truth).
\definecolor{pnlA}{RGB}{226,238,252}% pastel blue
\definecolor{pnlB}{RGB}{255,243,214}% pastel yellow
\definecolor{pnlC}{RGB}{225,245,228}% pastel green
\definecolor{pnlD}{RGB}{238,231,250}% pastel lavender
\definecolor{accOp}{RGB}{214,96,0}%   orange: changed operator
\definecolor{accOut}{RGB}{196,24,72}% crimson: contested outcome
\newcommand{\figicon}[2][0.66cm]{\includegraphics[height=#1]{figures/icons/#2.pdf}}
\newcommand{\figmark}[1]{\raisebox{-0.22em}{\includegraphics[height=0.34cm]{figures/icons/#1.pdf}}}
\newcommand{\op}[1]{{\setlength{\fboxsep}{0.7pt}\colorbox{accOp!20}{\textcolor{accOp!85!black}{\ttfamily\bfseries #1}}}}
\begin{tikzpicture}[
  x=1cm, y=1cm, font=\scriptsize,
  panel/.style={draw=black!65, rounded corners=5pt, line width=0.6pt},
  badge/.style={circle, draw=black!70, fill=white, inner sep=0.6pt, minimum size=9pt, font=\scriptsize\bfseries},
  ttl/.style={anchor=west, font=\scriptsize\bfseries},
  bubble/.style={rectangle callout, draw=black!65, fill=white, rounded corners=2.5pt, inner sep=3pt, align=left},
  thought/.style={draw=black!65, fill=white, rounded corners=6pt, inner sep=3pt, align=left},
  code/.style={draw=black!35, fill=white, rounded corners=1.5pt, inner sep=2.5pt, align=left, font=\ttfamily\scriptsize},
  lbl/.style={inner sep=0pt, anchor=south west},
]
\def\W{4.13}\def\G{0.16}\def\V{0.40}\def\H{2.64}
\def\Ix{0.52}% x-centre of the persona/action icon in every panel
% ---------------- panel frames ----------------
\path[panel, fill=pnlA] (0,\H+\V) rectangle ++(\W,\H);
\path[panel, fill=pnlB] (\W+\G,\H+\V) rectangle ++(\W,\H);
\path[panel, fill=pnlC] (0,0) rectangle ++(\W,\H);
\path[panel, fill=pnlD] (\W+\G,0) rectangle ++(\W,\H);
\foreach \n/\x/\y/\t in {1/0/\H+\V/A developer reads code, 2/\W+\G/\H+\V/A 14B model predicts,
                         3/0/0/Run both versions, 4/\W+\G/0/Affected-set claims exact?}{
  \node[badge] at (\x+0.24,\y+\H-0.24) {\n};
  \node[ttl] at (\x+0.44,\y+\H-0.24) {\t};
}
% ---------------- panel 1: the function and the question ----------------
\node[code, anchor=north east] (src) at (\W-0.14,2*\H+\V-0.48)
  {def f(l, t):\\[-1pt]
   \ \ return all(x \op{<} t\\[-1pt]
   \ \ \ \ \ \ \ \ \ \ \ \ \ for x in l)};
\node (dev) at (\Ix,\H+\V+0.50) {\figicon{technologist}};
\node[thought, anchor=west] (q) at (1.30,\H+\V+0.62) {What if \op{<} became \op{<=}?};
\fill[white, draw=black!65] (0.98,\H+\V+0.80) circle (0.055);
\fill[white, draw=black!65] (1.12,\H+\V+0.72) circle (0.075);
% ---------------- panel 2: the model's answers (real probe output, stripped variant) ----------------
\node[anchor=north west, inner sep=0pt] at (\W+\G+0.44,2*\H+\V-0.52) {Input \texttt{l=[1,3,4]}, \texttt{t=4}:};
\node (bot) at (\W+\G+\Ix,\H+\V+0.92) {\figicon{robot}};
\node[bubble, anchor=west, callout absolute pointer={([xshift=1pt]bot.east)}] (ans) at (\W+\G+1.20,\H+\V+0.92)
  {\begin{tabular}{@{}l@{\ \ }l@{\ \ }c@{}}
   With \op{<}:  & \texttt{False} & \figmark{check}\\[1pt]
   With \op{<=}: & \textcolor{accOut}{\underline{\texttt{\bfseries False}}} & \figmark{cross}
   \end{tabular}};
% ---------------- panel 3: executed contrast (witness) ----------------
\node[code, anchor=north west] (out) at (1.10,\H-0.78)
  {\begin{tabular}{@{}l@{ }l@{\ \ }c@{}}
   x \op{<}\ \ t $\to$ & False & \\[1pt]
   x \op{<=} t $\to$   & \textcolor{accOut}{\bfseries True} & \figmark{check}
   \end{tabular}};
\node (run) at (\Ix,0|-out.center) {\figicon{play}};
\node[anchor=north west, inner sep=0pt] at ([yshift=-7pt]out.south west) {Witness: \texttt{l=[1,3,4]}, \texttt{t=4}};
% ---------------- panel 4: checkable affected-input claims (ungrounded predicates) ----------------
\def\Mx{\W+\G+2.02}% column of the check/cross marks
\node[lbl, anchor=west] (la) at (\W+\G+1.10,\H-0.66) {Haiku:};
\node[inner sep=0pt] (ma) at (\Mx,0|-la.west) {\figmark{check}};
\node[anchor=west, inner sep=0pt] at ([xshift=3pt]ma.east) {exact on all tests};
\node[code, anchor=north west] (pa) at ([yshift=-4pt]la.south west)
  {any(x==t \ldots)\\[-1pt] and all(x\op{<=}t \ldots)};
\node[lbl, anchor=west] (lb) at ([yshift=-7pt]pa.south west) {14B:};
\node[inner sep=0pt] (mb) at (\Mx,0|-lb.west) {\figmark{cross}};
\node[anchor=west, inner sep=0pt] at ([xshift=3pt]mb.east) {not exact};
\node[code, anchor=north west] (pb) at ([yshift=-4pt]lb.south west) {any(x==t \ldots)};
\node (mag) at (\W+\G+\Ix,0|-pa.south) {\figicon{magnifier}};
% ---------------- arrow: guess (panel 2) -> truth (panel 3), labelled on the arrow ----------------
\draw[accOut, line width=0.9pt, -{Stealth[length=5pt]}, rounded corners=5pt]
  ([yshift=-1pt]ans.south) |- (\W+\G/2,\H+\V/2) -| ([xshift=-6pt]out.north east);
\node[fill=white, draw=accOut, rounded corners=2pt, inner sep=1.5pt, text=accOut, font=\scriptsize\bfseries]
  at (\W+\G/2,\H+\V/2) {run both};
\end{tikzpicture}

\caption{Counterfactual explanation as an executed contrast (real example: HumanEval+/52, \texttt{<}$\to$\texttt{<=};
display simplified). The 14B model mispredicts the changed output; running both versions gives the witness; affected-set
claims written as code are checked on all tests. Icons: Twemoji (CC-BY 4.0).}\label{fig:teaser}
\end{figure}

\noindent\fbox{\parbox{\dimexpr\columnwidth-2\fboxsep-2\fboxrule}{\small
\textbf{New idea:} explain a function by an \emph{executed contrast} whose claims (witness, affected-input set) are
themselves executable, and measure both their reliability and, next, their effect on readers' transfer to unseen changes.
\textbf{Why new:} counterfactual questions over code probe \emph{models}~\cite{vashishtha2026executable,hasanov2026path,yang2025reversion};
Cito et al.~\cite{cito2022counterfactual} explain \emph{classifiers} of code; WatChat~\cite{chandra2024watchat} and the
Whyline~\cite{ko2008whyline} are human-facing but neither checks the explanation against an executed contrast nor tests
transfer.
\textbf{Most related paper:} Vashishtha et al., \emph{Executable Counterfactuals}~\cite{vashishtha2026executable}.}}
\smallskip

\textbf{Hypotheses.}
\textbf{H0} (pipeline): without execution evidence, LLM explanations of an intervention's effect are often wrong;
grounding them in differential execution reduces the error.
\textbf{H1} (instrument): counterfactual (change-impact) questions are harder than matched descriptive ones.
\textbf{H2} (treatment): counterfactual explanations improve accuracy on \emph{unseen} change-impact questions over
code only and descriptive explanations. \textbf{H3} (calibration): they reduce overconfidence.

\textbf{Contributions.} (1) The primitive and \CoDE, a pipeline that intervenes, executes, articulates and checks
(\S\ref{sec:approach}). (2) H0 on \SNMutants{} HumanEval+/MBPP+ functions with two generators and two executable checks
(witness validity, affected-set exactness) (\S\ref{sec:eval}). (3) A model-side pilot of the H1 instrument
(LLMs as stand-in readers, \emph{not} human data). (4) A planned H2--H3 protocol
with generated held-out materials (\S\ref{sec:agenda}); no human study was run.
